\documentclass[12pt,a4paper]{article}

\usepackage[a4paper,margin=1in]{geometry}
\usepackage[T1]{fontenc}
\usepackage[utf8]{inputenc}
\usepackage{lmodern}
\usepackage{setspace}
\usepackage{enumitem}
\usepackage{array}
\usepackage{longtable}
\usepackage{booktabs}
\usepackage{tabularx}
\usepackage{titlesec}
\usepackage[hidelinks]{hyperref}
\usepackage{xcolor}
\usepackage{listings}
\usepackage{tikz}
\usepackage{graphicx}
\usetikzlibrary{positioning,shapes.geometric,arrows.meta}

\setstretch{1.12}
\setlength{\parindent}{0pt}
\setlength{\parskip}{6pt}
\titleformat{\section}{\normalfont\large\bfseries}{\thesection}{1em}{}
\titleformat{\subsection}{\normalfont\normalsize\bfseries}{\thesubsection}{1em}{}

\lstset{
  language=SQL,
  basicstyle=\ttfamily\scriptsize,
  frame=single,
  breaklines=true,
  showstringspaces=false,
  columns=fullflexible
}

\newcolumntype{Y}{>{\raggedright\arraybackslash}X}

\tikzset{
  entity/.style={draw, rectangle, rounded corners=2pt, text width=2.9cm, minimum height=0.9cm, align=center, fill=blue!5, font=\footnotesize},
  rel/.style={draw, diamond, aspect=2, text width=1.8cm, align=center, fill=orange!10, inner sep=1.2pt, font=\footnotesize},
  note/.style={font=\scriptsize, align=center},
  line/.style={-Latex, thick}
}

\begin{document}

\begin{center}

\includegraphics[width=0.22\textwidth]{Assets/logo.png}
\end{center}

\begin{center}
\textbf{UNIVERSITY OF CHITTAGONG}\\
Department of Computer Science and Engineering
\end{center}

\vspace{0.8cm}

\begin{center}
\textbf{\Large Report: Task 3}\\
\textbf{Refinement, Normalization, and SQL DDL for a Medical Database}
\end{center}

\vspace{0.8cm}

\begin{center}

\textbf{Course Title:} Database Management System\\
\textbf{Course Code:} CSE 414

\vspace{0.8cm}

\textbf{Submitted to:}\\
Dr. Rudra Pratap Deb Nath\\
Professor\\
Department of Computer Science and Engineering\\
University of Chittagong

\vspace{0.8cm}

\textbf{Submitted by:}\\
Name: Sanjida Mahi\\
ID: 24701009\\
Session: 2023--2024 B.Sc. Engineering (4th Semester)\\
Department of Computer Science and Engineering\\
University of Chittagong

\vspace{0.8cm}

\textbf{Date of Submission:}\\
03 May 2026

\end{center}

\newpage

\section{Introduction}

In Task 1, we built an initial relational schema and SQL DDL for the medical database. 
In Task 3, we're refining that original result logically by identifying functional
dependencies, checking normal forms, and adapting the schema where necessary.

We're using the supplied Excel workbook from \textbf{Tutorial 2} as the main source of data for this task. 
It is a denormalized source because one row may contain multiple side effects, multiple drug interactions,
and multiple clinical-trial conditions. 
Therefore, we're structuring the normalized schema to store these informations in separate relations.

\section{Relational Schema From Task 1}

\begin{lstlisting}
Drug(drug_name, drug_category)
Disease(disease_name, disease_category)
SideEffect(side_effect_name)
Researcher(main_researcher)
Company(company_name)
Product(product_name)
ClinicalTrial(
    clinical_trial_title,
    clinical_trial_status,
    clinical_trial_country,
    clinical_trial_department,
    clinical_trial_institution,
    clinical_trial_start_date,
    clinical_trial_completion_date,
    clinical_trial_condition,
    clinical_trial_participants
)
Studies(clinical_trial_title, drug_name)
Treats(drug_name, disease_name)
HasSideEffect(drug_name, side_effect_name)
InteractsWith(drug_name, interacts_with_drug_name)
ConductedBy(clinical_trial_title, main_researcher)
FocusesOn(clinical_trial_title, disease_name)
IsMarketedAs(product_name, drug_name)
ProducedBy(product_name, company_name)
\end{lstlisting}

\section{Functional Dependencies}
This section identifies the functional dependencies (FDs) that holds the
relations obtained from the Task 1 schema. These dependencies are later used as the basis for
normalization into Third Normal Form (3NF).

\newpage

\begin{lstlisting}
Drug(drug_name, drug_category)
    intended FD:
        drug_name -> drug_category
    but this is not safe for the workbook because one drug
    may belong to multiple categories

Disease(disease_name, disease_category)
    disease_name -> disease_category

SideEffect(side_effect_name)
    no non-trivial FD beyond the key

Researcher(main_researcher)
    no non-trivial FD beyond the key

Company(company_name)
    no non-trivial FD beyond the key

Product(product_name)
    no non-trivial FD inside this relation itself
    later joins suggest:
        product_name -> drug_name, company_name

ClinicalTrial(clinical_trial_title, ...)
    clinical_trial_title ->
        clinical_trial_status,
        clinical_trial_country,
        clinical_trial_department,
        clinical_trial_institution,
        clinical_trial_start_date,
        clinical_trial_completion_date,
        clinical_trial_condition,
        clinical_trial_participants
    clinical_trial_condition is not atomic in the workbook

Studies(clinical_trial_title, drug_name)
    no non-trivial FD beyond the composite key

Treats(drug_name, disease_name)
    no non-trivial FD beyond the composite key

HasSideEffect(drug_name, side_effect_name)
    no non-trivial FD beyond the composite key

InteractsWith(drug_name, interacts_with_drug_name)
    no non-trivial FD beyond the composite key

ConductedBy(clinical_trial_title, main_researcher)
    clinical_trial_title -> main_researcher
    main_researcher -> clinical_trial_title
    because Task 1 used UNIQUE on main_researcher

FocusesOn(clinical_trial_title, disease_name)
    clinical_trial_title -> disease_name
    this is too restrictive for the workbook

IsMarketedAs(product_name, drug_name)
    product_name -> drug_name

ProducedBy(product_name, company_name)
    product_name -> company_name
\end{lstlisting}

\subsection{Problems Revealed by the Task 1 Schema}
\begin{itemize}[leftmargin=1.5em]
    \item \textbf{Drug category was stored inside \texttt{Drug}.} This breaks down when one drug belongs to multiple categories in the sample data.
    \item \textbf{Clinical trial condition was stored as one attribute.} The workbook shows repeated trial-condition columns, so the attribute is multi-valued.
    \item \textbf{\texttt{FocusesOn} allowed only one disease per trial.} The workbook shows that a trial may be associated with several conditions.
    \item \textbf{Product-to-drug and product-to-company were split across two relations.} This is valid, but it can be simplified into one normalized \texttt{Product} relation without losing information.
    \item \textbf{Institution details were embedded in \texttt{ClinicalTrial}.} Separating them gives better reuse and avoids repeated location facts.
\end{itemize}

\section{Normalization to 3NF}
Starting from Task 1, the schema needed the following refinements:

\begin{enumerate}[leftmargin=1.5em]
    \item Replace the single category attribute in \texttt{Drug}.\\
    Introduce \texttt{DrugCategory} and \texttt{DrugHasCategory}.
    \item Replace the single trial-condition attribute in \texttt{ClinicalTrial}.\\
    Introduce \texttt{ConditionStudied} and \texttt{ClinicalTrialCondition}.
    \item Replace the overly restrictive \texttt{FocusesOn} design with a proper multi-valued trial-condition structure.
    \item Merge \texttt{Product}, \texttt{IsMarketedAs}, and \texttt{ProducedBy} into one cleaner \texttt{Product} relation.
    \item Extract \texttt{Institution} and \texttt{Researcher} so repeated trial-location facts are stored once.
\end{enumerate}

After these changes, the resulting schema became the normalized schema presented later in this
report.

\subsection{Relation-by-Relation Refinement}
\begin{enumerate}[leftmargin=1.5em]
    \item {\small\texttt{Drug(drug\_name, drug\_category)}}\\
    Decomposed into \texttt{Drug}, \texttt{DrugCategory}, and \texttt{DrugHasCategory} because one drug may belong to multiple categories.

    \item {\small\texttt{Disease(disease\_name, disease\_category)}}\\
    \item 
    Refined into \texttt{Disease} and \texttt{DiseaseCategory} to separate category descriptions cleanly.

    \item {\small\texttt{ClinicalTrial(..., clinical\_trial\_condition, ... )}}\\
    Decomposed into \texttt{ClinicalTrial}, \texttt{ConditionStudied}, and\\
    \texttt{ClinicalTrialCondition} because trial condition is multi-valued in the workbook.

    \item {\small\texttt{Product, IsMarketedAs, ProducedBy}}\\
    Combined logically into one normalized \texttt{Product(product\_name, company\_name, drug\_name)} relation plus separate \texttt{Company}.

    \item {\small\texttt{ConductedBy} and trial location attributes inside \texttt{ClinicalTrial}}\\
    Refined through separate \texttt{Institution} and \texttt{Researcher} relations so repeated location and researcher facts are stored once.

    \item {\small\texttt{FocusesOn}}\\
    Replaced conceptually by \texttt{ClinicalTrialCondition}, because the original one-trial-to-one-disease restriction does not fit the sample workbook.

    \item {\small\texttt{Studies, Treats, HasSideEffect, InteractsWith}}\\
    Kept as bridge relations because they already model many-to-many facts correctly.
\end{enumerate}

\subsection{3NF Check}
A relation is in 3NF if for every non-trivial FD \texttt{X $\rightarrow$ A}, at least one of the
following holds:

\begin{enumerate}[leftmargin=1.5em]
    \item \texttt{X} is a superkey, or
    \item \texttt{A} is a prime attribute.
\end{enumerate}

For the final schema produced from the Task 1 starting point:

\begin{itemize}[leftmargin=1.5em]
    \item In every entity relation with a non-key attribute, the determinant is the primary key.
    \item In every bridge relation, there are no non-key attributes at all.
    \item Therefore, all resulting relations satisfy 3NF directly.
\end{itemize}

\section{Highest Normal Form and BCNF Check}
The next step is to determine whether the relations also satisfy BCNF. A relation is in BCNF if,
for every non-trivial FD \texttt{X $\rightarrow$ A}, \texttt{X} is a superkey.

\begin{enumerate}[leftmargin=1.5em]
    \item {\small\texttt{Drug}}: BCNF. Only key attribute exists.
    \item {\small\texttt{DrugCategory}}: BCNF. Only key attribute exists.
    \item {\small\texttt{DrugHasCategory}}: BCNF. Composite key only; no non-key attributes.
    \item {\small\texttt{DiseaseCategory}}: BCNF. Only key attribute exists.
    \item {\small\texttt{Disease}}: BCNF. \texttt{disease\_name} is the key and determines the only non-key attribute.
    \item {\small\texttt{SideEffect}}: BCNF. Only key attribute exists.
    \item {\small\texttt{DrugTreats}}: BCNF. Composite key only; no non-key attributes.
    \item {\small\texttt{DrugHasSideEffect}}: BCNF. Composite key only; no non-key attributes.
    \item {\small\texttt{DrugInteraction}}: BCNF. Composite key only; no non-key attributes.
    \item {\small\texttt{Company}}: BCNF. Only key attribute exists.
    \item {\small\texttt{Product}}: BCNF. \texttt{(product\_name, company\_name)} is the determinant and a key.
    \item {\small\texttt{Institution}}: BCNF. \texttt{(institution\_name, country)} is the determinant and a key.
    \item {\small\texttt{Researcher}}: BCNF. No non-trivial FD beyond the composite key.
    \item {\small\texttt{ClinicalTrial}}: BCNF. \texttt{trial\_title} is the determinant and a key.
    \item {\small\texttt{ClinicalTrialStudiesDrug}}: BCNF. Composite key only; no non-key attributes.
    \item {\small\texttt{ConditionStudied}}: BCNF. Only key attribute exists.
    \item {\small\texttt{ClinicalTrialCondition}}: BCNF. Composite key only; no non-key attributes.
\end{enumerate}

Thus, the normalized schema derived from Task 1 is not only in 3NF but also in BCNF under the
stated assumptions.

\section{Normalized Relations and How They Were Derived}
\subsection{Derivation from Task 1 and the Spreadsheet}
The spreadsheet was effectively one wide denormalized relation containing:

\begin{itemize}[leftmargin=1.5em]
    \item drug facts,
    \item disease facts,
    \item repeated side-effect columns,
    \item repeated interacting-drug columns,
    \item product and company facts,
    \item repeated clinical-trial facts,
    \item repeated clinical-trial-condition columns.
\end{itemize}

The normalized relations were derived from the Task 1 schema by applying the following ideas:

\begin{enumerate}[leftmargin=1.5em]
    \item Separate multi-valued facts into bridge relations.
    \item Separate descriptive categories into their own tables.
    \item Keep one fact in one place only.
    \item Preserve natural identifiers where they are meaningful and sufficiently stable for the mini-world.
    \item Remove Task 1 restrictions that do not fit the sample workbook, especially single-category drugs and single-condition trials.
\end{enumerate}

\subsection{Final Normalized Relations}
\begin{lstlisting}
Drug(
    drug_name VARCHAR(120) PK
)

DrugCategory(
    category_name VARCHAR(120) PK
)

DrugHasCategory(
    drug_name VARCHAR(120) PK/FK -> Drug,
    category_name VARCHAR(120) PK/FK -> DrugCategory
)

DiseaseCategory(
    category_name VARCHAR(120) PK
)

Disease(
    disease_name VARCHAR(200) PK,
    category_name VARCHAR(120) FK -> DiseaseCategory
)

SideEffect(
    side_effect_name VARCHAR(120) PK
)

DrugTreats(
    drug_name VARCHAR(120) PK/FK -> Drug,
    disease_name VARCHAR(200) PK/FK -> Disease
)

DrugHasSideEffect(
    drug_name VARCHAR(120) PK/FK -> Drug,
    side_effect_name VARCHAR(120) PK/FK -> SideEffect
)

DrugInteraction(
    drug_name VARCHAR(120) PK/FK -> Drug,
    interacting_drug_name VARCHAR(120) PK/FK -> Drug
)

Company(
    company_name VARCHAR(160) PK
)

Product(
    product_name VARCHAR(160),
    company_name VARCHAR(160),
    drug_name VARCHAR(120),
    PK(product_name, company_name),
    FK company_name -> Company,
    FK drug_name -> Drug
)

Institution(
    institution_name VARCHAR(200),
    country VARCHAR(100),
    facility_address VARCHAR(255),
    PK(institution_name, country)
)

Researcher(
    researcher_name VARCHAR(160),
    institution_name VARCHAR(200),
    country VARCHAR(100),
    PK(researcher_name, institution_name, country),
    FK(institution_name, country) -> Institution
)

ClinicalTrial(
    trial_title VARCHAR(255) PK,
    start_date DATE,
    completion_date DATE,
    participants INT,
    status VARCHAR(80),
    institution_name VARCHAR(200),
    country VARCHAR(100),
    main_researcher_name VARCHAR(160)
)

ClinicalTrialStudiesDrug(
    trial_title VARCHAR(255) PK/FK -> ClinicalTrial,
    drug_name VARCHAR(120) PK/FK -> Drug
)

ConditionStudied(
    condition_name VARCHAR(200) PK
)

ClinicalTrialCondition(
    trial_title VARCHAR(255) PK/FK -> ClinicalTrial,
    condition_name VARCHAR(200) PK/FK -> ConditionStudied
)
\end{lstlisting}

Key derivation ideas:
\begin{itemize}[leftmargin=1.5em]
    \item \texttt{DrugHasCategory} was introduced because one drug may have multiple categories.
    \item \texttt{ClinicalTrialCondition} was introduced because one trial may mention multiple conditions.
    \item \texttt{Product} was kept as one normalized relation instead of splitting product-to-drug and product-to-company facts across separate bridge tables.
\end{itemize}

\section{SQL DDL}
The final SQL DDL is included below and also saved separately as:

\texttt{24701026\_CSE414\_Task3\_Normalized\_Schema.sql}

\lstinputlisting{24701026_CSE414_Task3_Normalized_Schema.sql}

\section{Revised ER Diagram}
The original Task 1 ER and relational design needed revision, but not a completely different
mini-world interpretation. Task 3 makes the following refinements explicit:

\begin{itemize}[leftmargin=1.5em]
    \item composite natural keys are retained where they help avoid unsafe uniqueness assumptions,
    \item bridge relations are justified by the observed multi-valued spreadsheet fields,
    \item candidate-key assumptions are turned into enforceable SQL constraints.
\end{itemize}

The revised ER diagram is therefore the corrected version of the original design and is close to the
stronger structure that later appeared in Task 2:

\begin{figure}[h]
    \centering
    \includegraphics[width=\textwidth]{Assets/ERD.jpg}
    \caption{Revised ER diagram for the normalized medical database.}
\end{figure}

\section{Comparison with Tasks 1 and 2}
\subsection{Compared with Task 1}
\begin{itemize}[leftmargin=1.5em]
    \item Task 1 stored \texttt{drug\_category} inside \texttt{Drug}; Task 3 replaces this with a many-to-many category design.
    \item Task 1 stored a single \texttt{clinical\_trial\_condition}; Task 3 models trial conditions as a separate multi-valued relation.
    \item Task 1 used \texttt{FocusesOn} with one disease per trial; Task 3 removes that restriction.
    \item Task 3 is formally justified through functional dependencies and normal-form analysis.
\end{itemize}

\subsection{Compared with Task 2}
\begin{itemize}[leftmargin=1.5em]
    \item The final Task 3 structure is close to Task 2 because Task 2 had already corrected most conceptual flaws from Task 1.
    \item The difference is that Task 3 now derives that structure from the original Task 1 schema and justifies it with dependencies and normal forms.
    \item Task 3 also turns the design into explicit SQL constraints and a normalization argument.
\end{itemize}

\subsection{Advantages of the Task 3 Design}
\begin{itemize}[leftmargin=1.5em]
    \item Less redundancy than the original spreadsheet.
    \item Lower risk of update anomalies, insertion anomalies, and deletion anomalies.
    \item Clearer enforcement of keys and referential integrity.
    \item Better readiness for later querying and data loading.
\end{itemize}

\section{Conclusion}
The Task 1 medical database schema required refinement because some of its assumptions did not fit
the sample workbook, especially in the handling of drug categories and clinical-trial conditions.
After identifying the functional dependencies and revising the schema, the final result was found to
satisfy 3NF and, under the stated assumptions, also BCNF. The final SQL DDL therefore preserves a
clean, non-redundant structure that is suitable for the next project stages.

\end{document}
